\documentclass[10pt,a4paper]{amsart}
\usepackage[margin=19mm]{geometry}
\usepackage{amsmath,amssymb,mathtools}
\usepackage[scr=rsfs,cal=euler]{mathalfa}
\usepackage{xfrac}
\usepackage[unicode,hidelinks,bookmarks=false]{hyperref}
\newcommand{\ie}{i.e.\ }
\newcommand{\cf}{cf.\ }
\newcommand{\resp}{respectively\ }
\newcommand{\Subsection}{\subsection}
\providecommand{\nobreakdash}{\nobreak\hbox{-}}
\setlength{\emergencystretch}{2em}
\multlinegap0pt
\newtheorem{thm}{Theorem}[section]
\newtheorem{prop}[thm]{Proposition}
\newtheorem{cor}[thm]{Corollary}
\newtheorem{lem}[thm]{Lemma}
\newtheorem*{thm1}{Theorem 1}
\newtheorem*{thm2}{Theorem 2}
\newcommand{\bF}{\mathbf{F}}
\begin{document}
\title[The space of two-dimensional vectors over a four-dimensional]{The space of two-dimensional vectors over a four-dimensional division algebra over $F_2$}
\author{Daisuke Tambara}
\date{}
\maketitle
\noindent\emph{Source and licence.} The space of two-dimensional vectors over a four-dimensional division algebra over $F_2$. Version arXiv:2607.20932v1 (CC BY 4.0). This material is distributed under the Creative Commons Attribution 4.0 International licence, http://creativecommons.org/licenses/by/4.0/, as recorded in the arXiv OAI licence metadata for this exact version and shown on the abstract page https://arxiv.org/abs/2607.20932v1. Full rights notice: licenses/arxiv-2607.20932-CC-BY-4.0.txt.\par\smallskip
\noindent\textit{Editorial scope.} Counted unit: the original \texttt{thm} environment \texttt{-} at source lines 1975--1984 together with its complete proof at lines 1987--2056; the delivered document keeps source lines 1906--2057 so that every referenced label and every citation resolves inside it. Native editable source is the paper's own concatenated TeX; the AMS wrapper is editorial formatting only. No publisher class, upstream script or unrelated artwork is redistributed.
\section{The algebra of type W}

In this section we prove Theorem 2.2.

Let $E/F$ be a quadratic extension with nontrivial automorphism $\sigma$.
Let $\omega\in E-F$.
Let $A$ be the algebra of type W associated with $E/F$:
$A=E\oplus E$, 
$(a,b)(c,d)=(ac+\omega b^{\sigma}d, bc+a^{\sigma}d)$.
This is a division algebra,  four-dimensional with basis
$(1,0), (\omega,0), (0,1), (0,\omega)$.
Put $e=(1,0)$, $j=(0,1)$. $e$ is the identity element.

Put $A_0=E\oplus 0=\{(a,0)\mid a\in E\}$.
This is a subalgebra isomorphic to the field $E$.
The multiplication rule specializes to
\begin{align*}
(a,0)(c,d)&=(ac, a^{\sigma}d),\\
(a,b)(c,0)&=(ac,bc).
\end{align*}
The left and right multiplication by elements of $A_0$  is an associative action 
on  $A$, so that $A$ becomes a two-sided $A_0$-module.
And
$(a,b)=(a,0)+j(b,0)$,
$A=A_0+jA_0$.

Moreover one sees 
$$
t(xy)=(tx)y,\; (xt)y=x(ty),\; (xy)t=x(yt)
$$
for $x,y\in A$, $t\in A_0$ (\cite{Knu}).


A computation shows

\begin{lem}
For any $x\in A$ we have $(jx)j=j(xj)$ if and only if $x\in A_0$.
\end{lem}



\begin{prop}
If $0\ne t\in A_0$, then $A(tx,ty)=A(x,y)$ for any $x,y\in A$.

\end{prop}

\begin{proof}
By the middle associativity for $t\in A_0$
$$
a(tx,ty)=(a(tx),a(ty))=((at)x,(at)y)=(at)(x,y).
$$
Hence
$A(tx,ty)=(At)(x,y)=A(x,y)$
for $0\ne t\in A_0$.
\end{proof}


\begin{prop}
If $T\in GL_2(A_0)$, then
the linear transformation $(x,y)\mapsto (x,y)T$ of $A^2$ maps the space
$A(x,y)$ isomorphically to $A((x,y)T)$.
\end{prop}

\begin{proof}
Follows  from the right associativity $(ax)t=a(xt)$ for $t\in A_0$.
\end{proof}

We restate Theorem 2.2 through the isomorphism $E\cong A_0$.


\begin{thm}
Let $x,y,x',y'\in A$ be nonzero elements.
Then $A(x,y)=A(x',y')$ if and only if either of the following holds:

(i) $x'=\lambda x$, $y'=\lambda y$ for some $\lambda\in A_0^{\times}$.

(ii) $x=y\mu$, $x'=y'\mu$
 for some $\mu\in A_0^{\times}$.
 
\end{thm}

\begin{proof}
(i) is sufficient for the equality $A(x,y)=A(x',y')$ by Proposition 6.2.

Assume (ii):
$x=y\mu$, $x'=y\mu'$ with $\mu\in A_0^{\times}$.
Put
$$
T=\begin{pmatrix} 1 & 0 \\ -\mu & 1 \end{pmatrix},
$$
so that
$(x,y)T=(0,y)$, $(x',y')T=(0,y')$.
Since $A$ is a division algebra, we have $A(0,y)=0\oplus A=A(0,y')$.
By Proposition 6.3 it follows that $A(x,y)=A(x',y')$.




Let us prove the necessity.
Write $x=(a,b)$, $y=(c,d)$ and
put
$$
T=\begin{pmatrix} (a,0) & (c,0) \\
(b,0) & (d,0) \end{pmatrix}
\in M_2(A_0),
$$
so that
$(x,y)=(e,j)T$.


First consider the case where $x,y$ are right linearly independent over $A_0$.
Then $T^{-1}$ exists and
$(x,y)T^{-1}=(e,j)$.
Suppose $A(x,y)=A(x',y')$.
Put $(x',y')T^{-1}=(\tilde x',\tilde y')$.
Then, by Proposition 6.3,
$A(e,j)=A(\tilde x',\tilde y')$.
Since
$$
A(e,j)=\{(ae,aj)\mid a\in A\}=\{(a,aj)\mid a\in A\},
$$
the last equation  means that
$(a\tilde x')j=a\tilde y' $ for all $a\in A$.
Letting $a=e, j$, we have equations
$\tilde x'j=\tilde y'$, 
$(j\tilde x')j=j\tilde y'$, 
hence
$(j\tilde x')j=j(\tilde x' j)$.
By Lemma 6.1 this means   $\tilde x'\in A_0$.
Put $\lambda=\tilde x'$. Then
$(\tilde x',\tilde y')=(\lambda, \lambda j)=\lambda(e,j)$, so
$$
(x',y')=(\lambda(e,j))T=\lambda((e,j)T)=\lambda(x,y).
$$
Thus (i) holds.


Next consider the case where $x,y$ are right linearly dependent over $A_0$.
 Write $y=xt$ with $t\in A_0$.
For any $a\in A$
$$
a(x,y)=(ax,ay)=(ax,a(xt))=(ax,(ax)t),
$$
hence
$A(x,y)=\{(a,at)\mid a\in A\}$.
Therefore,
if $A(x,y)=A(x',y')$, then $y'=x't$, so (ii) holds.



\end{proof}


\begin{thebibliography}{99}
\bibitem[Knu]{Knu} D.E.Knuth, Finite semifields and projective planes, J.\ Algebra \textbf{2} (1965), 182--217.
\end{thebibliography}
\end{document}
